\chapter{Sequential control inventory}
\label{ch:fsms}

This chapter lists every piece of sequential control in the delivery and states
plainly which are coded state machines and which are flag-encoded behaviour. The
distinction matters when reading the source: only one module has a state
register.

\begin{table}[H]
\centering
\caption{Sequential control across the delivery}
\label{tab:fsminv}
\begin{tabular}{@{}L{0.24\linewidth}L{0.14\linewidth}L{0.40\linewidth}c@{}}
\toprule
\textbf{Module} & \textbf{Kind} & \textbf{State elements} & \textbf{Figure} \\
\midrule
\file{lsu.v} & Coded FSM & \sig{state\_q[1:0]}: \sig{S\_IDLE}, \sig{S\_RD}, \sig{S\_WR},
  plus three handshake trackers & \ref{fig:lsufsm} \\
\file{fetch\_unit.v} & Flag-encoded & \sig{ar\_pending\_q}, \sig{inflight\_q},
  \sig{discard\_q}, \sig{hold\_valid\_q} & \ref{fig:fetchfsm} \\
\file{axi\_lite\_slave.v} & Flag-encoded & \sig{aw\_got\_q}, \sig{w\_got\_q},
  \sig{bvalid}, \sig{rvalid} & \ref{fig:slavewr} \\
\file{pic.v} (per source) & Flag-encoded & \sig{edge\_pend}, \sig{sw\_pend},
  \sig{active}, \sig{escalated}, \sig{spurious} & \ref{fig:picsrcfsm} \\
\file{pic.v} (nesting) & Counter + stack & \sig{depth[4:0]}, \sig{stack\_id},
  \sig{stack\_key} & \ref{fig:picdepth} \\
\file{cpu\_top.v} & Flag-encoded & \sig{in\_irq}, \sig{cpu\_in\_trap} & \ref{fig:trapfsm} \\
\file{branch\_predictor.v} & Table + pointer & BTB/BHT array, RAS pointer
  \sig{sp\_q} and count \sig{cnt\_q} & --- \\
\file{mtimer.v} & Counters only & \sig{mtime\_q}, \sig{mtimecmp\_q}, \sig{irq} & --- \\
\bottomrule
\end{tabular}
\end{table}

\begin{sloppypar}
The remaining modules --- \file{decode.v}, \file{imm\_gen.v}, \file{control.v},
\file{alu.v}, \file{alu\_top.v}, \file{branch\_unit.v},
\file{writeback\_mux.v}, \file{exception\_unit.v} and \file{hazard\_unit.v} ---
are purely combinational and contain no state at all.
\end{sloppypar}
\noindent In particular the hazard unit,
which in many designs is a state machine, is five combinational equations here.

\section{CPU trap state}

Two flags in \file{cpu\_top.v} track where the core is with respect to a trap.
They are separate because they answer different questions.

\begin{table}[H]
\centering
\caption{Trap-tracking flags}
\label{tab:trapflags}
\begin{tabular}{@{}lL{0.66\linewidth}@{}}
\toprule
\textbf{Flag} & \textbf{Purpose} \\
\midrule
\sig{in\_irq}      & Set when an \emph{interrupt} is taken, cleared on \sig{MRET}. Gates
  \sig{cpu\_irq\_eoi}, so an \sig{MRET} closing a synchronous trap does not pop a PIC
  nesting level. \\
\sig{cpu\_in\_trap} & Set on \emph{any} trap entry, cleared on \sig{MRET}. A nested
  synchronous trap inside a handler keeps it high. Exported for observability and
  assertions only; the PIC does not consume it. \\
\bottomrule
\end{tabular}
\end{table}

\begin{figure}[H]
\centering
\begin{tikzpicture}[x=1mm,y=1mm,>={Stealth[length=2.4mm]},semithick]
  \node[stbi] (run) at (0,0)    {RUNNING};
  \node[stb]  (exc) at (78,20)  {IN\_TRAP};
  \node[stb]  (irq) at (78,-24) {IN\_TRAP + IN\_IRQ};

  % synchronous exception leg
  \draw[->,draw=accent] (run) to[bend left=13]
    node[lbl,pos=0.42,above left=0mm and -6mm,align=center]
    {synchronous exception\\
     \scriptsize trap\_take \&\ !irq\_take \&\ s2\_advance} (exc);
  \draw[->,draw=accent] (exc) to[bend left=13]
    node[lbl,pos=0.35,below right=0mm and -5mm,align=center]
    {MRET\\\scriptsize no eoi pulse} (run);

  % interrupt leg
  \draw[->,draw=accent] (run) to[bend right=13]
    node[lbl,pos=0.42,below left=0mm and -6mm,align=center]
    {interrupt taken\\\scriptsize irq\_take \&\ s2\_advance} (irq);
  \draw[->,draw=accent] (irq) to[bend right=13]
    node[lbl,pos=0.35,above right=0mm and -5mm,align=center]
    {MRET\\\scriptsize emits cpu\_irq\_eoi\_o} (run);

  \draw[->,draw=accent,loop right,looseness=7,out=35,in=-35] (exc) to
    node[lbl,right=1mm,align=left]{nested synchronous\\\scriptsize trap: stays high} (exc);
  \draw[->,draw=accent,loop right,looseness=7,out=35,in=-35] (irq) to
    node[lbl,right=1mm,align=left]{nested synchronous\\\scriptsize trap: both stay high} (irq);
\end{tikzpicture}
\caption{CPU trap tracking. Only \sig{in\_irq} gates the end-of-interrupt
pulse, so exactly one is emitted per interrupt and none for an exception. Both
flags are cleared by the same \sig{MRET}, and a synchronous trap taken inside
either handler leaves them where they are.}
\label{fig:trapfsm}
\end{figure}
